\subsection{A submitted effect table now returns an action score}
\label{app:ioi-decision-replay}

An outside researcher should not have to infer decision quality from an MAE
table. The IOI decision replay takes the existing effect CSV and ObserverCard,
chooses actions using the submitted means, and returns their held-out losses.
It uses only stored measurements, so neither GPT-2 nor the submitted observer
implementation runs during scoring. The supplied example passes the published
attribution-patching predictions through exactly this boundary.

\paragraph{The frozen decision.}
The task version is
\path{ioi-gpt2-small-action-selection@phase5-noop-v1-b160}.
It uses the original Phase 5 panel: 160 calibration masks averaged over
192 training prompts, 320 candidate masks in ten pools, and all 81,920
candidate responses on 256 held-out prompts. Each pool also includes exact
no-op, whose response is analytically zero. For each of the existing targets
$t\in\{0.5,1,1.5\}$, the selector minimizes predicted mean distance to the
target, breaking ties by fewer heads and then mask ID. Selection never reads
held-out responses. The same mask then applies to every test prompt in that
pool. The score averages $|Y_i(m)-t|$ over prompts and pools.

\paragraph{What the researcher receives.}
The scorecard reports prediction MAE, selected masks, action loss at all three
targets, excess loss over no-op, and paired comparisons with raw AtP,
scalar-calibrated AtP, additive ridge, and all-pairs ridge. It also records
task and prediction hashes, scoring-code hashes, and the software environment.
Paired 95\% percentile intervals use 2,000 resamples of 107 unordered-name-pair
clusters and ten pools, with seed 20260904. Every prompt in a sampled cluster
stays together. Fits and selected actions stay fixed; these intervals do not
describe uncertainty from refitting an observer.

\begin{table}[H]
\centering
\small
\begin{tabular}{@{}lrrrr@{}}
\toprule
Observer & Prediction MAE & $R_{0.5}$ & $R_1$ & $R_{1.5}$ \\
\midrule
Raw AtP & 0.662 & 1.075 & 0.947 & 1.093 \\
Scalar-calibrated AtP & 0.632 & 0.972 & 0.999 & 1.114 \\
Additive ridge & 0.486 & 1.053 & 0.972 & 1.024 \\
All-pairs ridge & 0.395 & 1.018 & 1.022 & 1.006 \\
Exact no-op & --- & 0.500 & 1.000 & 1.500 \\
\bottomrule
\end{tabular}
\caption{The same submitted tables scored for prediction and action.
Lower is better. AtP uses white-box gradients; the fitted ridge observers use
finite forward measurements. Cached scoring does not imply equal construction
cost. This secondary open replay is separate from the clean-task confirmation.}
\label{tab:ioi-decision-replay}
\end{table}

The point estimates rank prediction and action differently, but do not
establish a reliable observer ordering at $t=1$. For calibrated minus raw AtP,
the action-loss difference is $0.053$, with interval $[-0.076,0.203]$.
None of the four observers establishes a gain over no-op at that target.
All four lose to no-op at $t=0.5$ and beat it at $t=1.5$ under their paired
intervals. A researcher therefore leaves with a decision comparison,
including evidence about whether to act, rather than only a prediction rank.

\paragraph{Replay, not another confirmation.}
This task was assembled after the original results were known and adds no new
model measurements. The responses are public so users can obtain local scores;
those scores are not sealed submissions. The original Phase 5 experiment and
the later Phase 7 clean-task confirmation keep their own identities. Training
on this replay's scoring responses would invalidate a held-out claim, even
though the software could still compute a number.
